\documentclass{article}
\usepackage[T1]{fontenc}
\usepackage{lmodern}
\usepackage{graphicx}
\usepackage{amsmath,amssymb}
\usepackage[a4paper,margin=2cm]{geometry}
\usepackage[absolute,overlay]{textpos}
\setlength{\TPHorizModule}{1mm}
\setlength{\TPVertModule}{1mm}
\usepackage[indonesian]{babel}
\usepackage{hyperref}
\setlength{\emergencystretch}{2em}
% unit: Halaman 14

\begin{document}

% === Halaman 14 (llm) ===

\section*{Medan Galois (Galois Field)}

\begin{itemize}
    \item Medan Galois adalah medan berhingga dengan $p^m$ elemen, $p$ adalah bilangan prima dan $m \ge 1$.
    \item Notasi: $GF(p^m)$
    \item Kasus paling sederhana: bila $m = 1 \rightarrow GF(p)$ atau $\text{F}_p$
    \item $GF(p)$ terdiri dari himpunan $\{0, 1, 2, ..., p-1\}$ dan operasi penjumlahan dan perkalian dilakukan di dalam modulus $p$.
\end{itemize}

\begin{enumerate}
    \item \textbf{Penjumlahan}, dilakukan dalam modulus $p$
    
    Jika $a, b \in GF(p)$, maka $a + b = r$, yang dalam hal ini $r = (a + b) \pmod p$, $0 \le r \le p - 1$
    
    \item \textbf{Perkalian}, dilakukan dalam modulus $p$
    
    Jika $a, b \in GF(p)$, maka $a \cdot b = s$, yang dalam hal ini $s = (a \cdot b) \pmod p$, $0 \le s \le p - 1$
\end{enumerate}

\vspace{10mm}

\centering Rinaldi Munir/II4021 Kriptografi

\end{document}
